\section{Compatible tie partners and their exact weighted mean}
\label{sec:tie-partners}

\status{The actual simultaneous-partner construction and the source's
dimension-weighted multiple-tie mean are proved here. This uses GG's
quantitative prefix projections, integral highest-vector generation and
restricted graded Gram estimate. The sharper gamma integrality constant
is proved explicitly. Segment and digit-class transfers and the actual
edge lifts are separate steps.}

\paragraph{TP1--TP3: the exact spaces and the sharper one-tie estimate.}
Work over $\mathcal O=\mathbb Z_p$, $K=\mathbb Q_p$, with $p\ne2$ and
$r\ge2$. A signed word has $N$ letters of fixed degrees. Let $Y$ be its
polynomial space subject to the actual prefix conditions
$\lambda_1\le M_k$ at lengths $J_k$, and put
$Y_{\mathcal O}=Y\cap\mathcal O[X]$. Its form is
$B_L=\mathrm{Fock}/\prod_jd_j!$.
A tie $i$ joins vector $a_i$ to covector $a_i+1$, whose degree is $n_i$;
the positions increase, the pairs are disjoint, and $z_i\in\mathbb N$.
Write $S_i$ for the preceding total degree and require $n_i\ge S_i$.

Before that covector, let $M_{i,w}$ be the space of full torus weight $w$
that is highest for the coordinates other than $\rho=r-1$, with all
cutoffs $J_k\le a_i+1$. TI3's explicit freeze map identifies
$M_{i,w}\cap\mathcal O[X]$ with the entire integral highest-vector lattice
of weight $w-n_i e_\rho$ after adjoining the covector. Indeed, on
$y_\rho\ne0$, its substitution is a common unitriangular change of vectors
and inverse transpose change of covectors. Highest-vector invariance
recovers a polynomial from its specialization $y=e_\rho$, which is the
integral inverse. This also preserves the old prefix conditions.
The order filtration in $x_{a_i,\rho}$ is saturated; lifting bases of its
successive coefficient images gives an integral adapted basis.

For a prefix of $K+1$ old letters, total $S$ and last vector degree $o$,
the exact strengthening of TI4 is
\[
 \begin{gathered}
 p^{c_I}\Gamma_{f,g}(n)=\sum_{j=0}^{S+o}c_j\binom nj,
 \qquad c_j\in\mathcal O,\\
 c_I=((K+4)r+6)\lfloor\log_p(S+o+r)\rfloor.
 \end{gathered}
 \tag{TP.1}
\]
Here $n\ge S$ and $z\in\mathbb N$. A short proof expands the contraction
as $D_\rho+D_\perp$. A paired monomial term becomes an integral-binomial
polynomial times a scalar. The two differentiated coefficients and the
$K+2$ row-factorial groupings lose at most $(K+4)rB$, where
$B=\lfloor\log_p(S+o+r)\rfloor$. The remaining factorials cancel exactly to
\[
 -v_p\binom S{j-a}-2v_p\binom oj-v_p\binom j{a'},
 \qquad 0\le a,a'\le j\le o.
\]
This loses at most $4B$, within the source's $6B$ allowance. The expanded
companion gives every coefficient and the binomial translation.

Let $b$ be the adapted-basis size and suppose the equal-order initial Gram
determinants have valuation at most $\kappa_4$ times their sizes, as GG
proves. For $A(n)=p^{\lceil\kappa\rceil}\Gamma(n)$, TI3 gives degree
$\sum_u(S+o-e_u)$ and leading coefficient
\[
 \prod_e\left[
 \left(\frac{p^{\lceil\kappa\rceil}e!\binom{z+e}e}{o!S!}\right)^{|I_e|}
 \det G_e\right].
\]
Only equal-order permutations attain the top degree. If
$2S,z+o,Y,r\le X$,
$\ell\ge\lfloor\log_pX\rfloor+1$ and
$\kappa\ge((K+4)r+6)\ell$, (TP.1) gives integral binomial coefficients.
Legendre's formula, $\lceil\kappa\rceil<\kappa+1$ and
$v_p\binom{z+e}e\le\ell-1$ give the precise leading-coefficient budget
$b(\kappa+3\ell+\kappa_4)$. TI5, with the integer
$h=\lfloor\log_pX\rfloor+1$, therefore yields
\[
 \underset{0\le j<Y}{\operatorname{Avg}}
 \#\{\text{Smith factors of }A(s+dj):v_p>E\}
 \le b\left[\frac{2S}{Y}
 +\frac{\kappa+3\ell+\kappa_4+2S\ell/Y}{E}\right],
 \tag{TP.2}
\]
for $E\ge1$, $s\ge S$, $d>0$ and $p\nmid d$.

\paragraph{TP4--TP7: one system of integral partners for all ties.}
Use GG's common projection and highest-vector-generation exponent $e_1$.
For one tie, Smith bases give integral highest vectors $h_x,h''_x$ and
scalars $d_x$ such that
$c_*\mathrm{Fock}(\phi_i h_x,h''_y)=\delta_{xy}d_x$ within each weight.
The normalization is
\[
 c_*=
 \frac{p^{\lceil\kappa\rceil}\binom{o+n_i+r+z_i-1}o}
 {(n_i-S)!S!\prod_{j\le a_i}d_j!}.
\]
Row-factorial grouping also makes $c_*\mathrm{Fock}(h_x,h_y)$ integral.
Call $x$ good when $v_p(d_x)\le E$, and set $e=\lfloor E\rfloor$.
Define on highest vectors
\[
 \varpi h_x=1_{x\,\mathrm{good}}h_x,\qquad
 \tau h_y=
 \sum_{\substack{x\,\mathrm{good}\\\lambda_x=\lambda_y}}
 \frac{p^e}{d_x}\,c_*\mathrm{Fock}(h_x,h_y)h''_x.
 \tag{TP.3}
\]
These coefficients are integral. Extend the maps equivariantly over $K$,
then coefficientwise over later variables. The projection onto the earlier
cutoffs costs $(a_i+2)e_1$; the pure-type projection and highest-vector
generation cost another $2e_1$. Thus the resulting $T_i,\pi_i$ satisfy
\[
 \begin{gathered}
 p^{K_i}T_i(Y_{\mathcal O})\subseteq Y_{\mathcal O},
 \qquad K_i=e_1(a_i+4)\le e_1(N+2),\\
 B_L(\phi_i u,T_i v)=p^eB_L(\pi_i u,v),
 \qquad \pi_i\phi_j=\phi_j\pi_i\quad(j>i).
 \end{gathered}
 \tag{TP.4}
\]
The pairing identity permits $u$ with only the cutoffs
$J_k\le a_i+1$; $v$ satisfies all cutoffs. It follows on highest vectors
from (TP.3), then on their cyclic modules by Fock contravariance.

Let $D_{i,w}(n)$ be the dimension of the actual continuation space
\texttt{YYsp}: multiply a nonzero frozen highest vector's cyclic module by
later polynomials and impose all cutoffs. If $b_{i,w}(n)$ is its bad Smith
count, the saturated lattice
\[
 L_1=Y_{\mathcal O}\cap\bigcap_i\ker(\pi_i-1)
 \quad\text{satisfies}\quad
 \operatorname{codim}L_1\le\sum_{i,w}b_{i,w}(n)D_{i,w}(n).
 \tag{TP.5}
\]
A primitive $u\in L_1$ is primitive in $Y_{\mathcal O}$ and has a unit
coefficient. Project its corresponding monomial onto all prefix cutoffs
and multiply by $p^{e_1(N+1)}$. This gives $z\in Y_{\mathcal O}$ with
$|B_L(u,z)|\ge p^{-e_1(N+1)}$: the normalized monomial pairing is the
inverse of a product of integer multinomials. Define
\[
 u'=(p^{K_0}T_0)\cdots(p^{K_{t-1}}T_{t-1})z,
 \qquad D=\lfloor t(E+\kappa)+\kappa\rfloor.
\]
If $e_1(N+2)\le\kappa$, peeling (TP.4) in this order gives
$|B_L((\phi_0\cdots\phi_{t-1})u,u')|\ge p^{-D}$.
Later ties preserve the needed earlier cutoffs; the partial inverses need
not commute. Smith form now supplies $\operatorname{rank}L_1$ integral
pairs with matrix $p^D I$. Thus the source's \texttt{bad} count is bounded
by (TP.5).

\paragraph{TP8--TP11: retaining the multiplicities in the average.}
Write $m_{i,w}=\dim M_{i,w}$. Rational highest-weight decomposition and the
freeze isomorphism give the exact identity
\[
 \sum_w m_{i,w}D_{i,w}(n)=\dim Y(n).
 \tag{TP.6}
\]
Each rational cyclic copy is preserved by the later-prefix projectors;
copies of the same highest weight have isomorphic continuation spaces.
No integral direct-sum assertion enters this identity.

Also $D_{i,w}$ increases with $n_i$. For an increase $u\ge0$, the frozen
highest vector changes from $h$ to $y_\rho^u h$. Project multiplication by
$y_\rho^u$ onto this Cartan component. The resulting map commutes with
raising operators, shifts weights and is nonzero on the highest line,
so its kernel is zero. It preserves later cutoffs because covector weights
are nonpositive and the Cartan projector commutes with the larger-prefix
action. Earlier cutoffs follow from freezing. This proves an injection of
the actual continuation spaces.

On the grid $n_i=s_i+d_i j$, $0\le j<Y_i$, let
$n_i^+=s_i+d_i(Y_i-1)$ and $S_i^+=S_i(n^+)\le s_i$.
Apply (TP.2) holding the other coordinates fixed, replace $D_{i,w}(n)$
by its endpoint value, and use (TP.6). Finite summation gives
\[
 \boxed{\begin{gathered}
 \operatorname{Avg}_n\operatorname{bad}(n;D)
 \le\sum_i\epsilon_i\operatorname{Avg}_n\dim Y(n^{(i,+)}),\\
 \epsilon_i=\frac{2S_i^+}{Y_i}
 +\frac{2\kappa+2S_i^+\ell/Y_i}{E}.
 \end{gathered}}
 \tag{TP.7}
\]
Here $n^{(i,+)}$ changes only coordinate $i$ to its endpoint. The general
conditions are
$\kappa\ge\tfrac32(N+2)r^2\ell$, $e_1(N+2)\le\kappa$ and
$3\ell+\kappa_4\le\kappa$, with the degree bounds used above.
In the literal alternating chain $N=2s+1$, the source size inequalities
$\operatorname{total}\le8a+r$,
$o_i+n_i^++r+z_i+1\le8a+r$, $Y_i\le8a+r$ and
$2S_i^++1\le8a+r$ allow exactly
\[
 \begin{gathered}
 \ell=\log_p(8a+r)+1,\quad e_1=\lfloor r^2\ell\rfloor,
 \\\kappa_4=(s+4)r^2\ell,\quad
 \kappa=3(s+2)r^2\ell.
 \end{gathered}
 \tag{TP.8}
\]
GG supplies the projection, highest-vector and Gram bounds at these constants.
Equations (TP.7)--(TP.8), with $a,E\ge1$
and the literal admissible cutoffs, are precisely \texttt{TheoremMV7}.
They retain the endpoint dimensions with their actual multiplicities.
\textup{[TP1--TP11]}
